\chapter{A Realidade Como É}

\lettrine{U}{ma reflexão inteligente é:} ``É assim que as coisas são''. O~Venerável
Buddhadāsa Bhikkhu, o renomado sábio Tailandês, disse: ``Se houvesse uma
inscrição útil para se colocar num medalhão à volta do pescoço, seria `É assim
que as coisas são'.'' Esta reflexão ajuda-nos a contemplar: onde quer que
estejamos, em qualquer momento e lugar, bom ou mau, ``É assim que as coisas
são''. É uma forma de trazer uma aceitação para dentro das nossas mentes, um
reparar em vez de uma reacção.

A prática da meditação é reflectir sobre ``como as coisas são'', a fim de ver os
medos e desejos que criamos. Esta é uma prática bastante simples, mas a prática
do \emph{Dhamma} devia ser muito, muito simples em vez de complicada. Muitos
métodos de meditação são muito, muito complicados, com muitos estágios e
técnicas -- então a pessoa torna-se viciada em coisas complicadas. Às vezes, por
causa do nosso apego às convicções, não sabemos realmente como são as coisas. No
entanto, quanto mais simples ficamos, mais claro, profundo e significativo tudo
é para nós.

Por exemplo, consideremos as pessoas aqui, os monges e monjas com quem vivemos.
Talvez nos sintamos atraídos por algumas pessoas, por outras talvez sintamos
aversão, com algumas simpatizamos, outras entendemos, algumas não entendemos;
mas qualquer que seja a ideia que tenhamos, podemos vê-la apenas como uma
``ideia'' de uma pessoa, ao invés de uma pessoa real. Podemos ouvir-nos a dizer:
``Não quero que ele seja assim\ldots{} Quero que ele seja diferente. Ele devia
ser desta outra forma\ldots{} não devia de ser assim.''

``Eu quero que seja diferente'' é o lamento dos tempos actuais, não é? Porque é
que a vida não pode ser diferente? Porque é que as pessoas têm que morrer?
Porque é que temos que envelhecer? Porquê esta doença? Porque é que temos que
estar separados dos nossos entes queridos? Porque é que crianças inocentes que
não magoariam ninguém nas suas vidas, idosos que não magoariam ninguém -- porque
é que eles têm que sofrer de fome ou de violência?

Há sempre alguma coisa nova e horrível a acontecer. No outro dia, alguém me
escreveu sobre os muçulmanos de Bangladesh a tentarem livrar-se das Tribos das
Colinas Budistas nas Colinas de Chatigão por genocídio. Depois ouvimos sobre os
Iranianos a tentarem erradicar os Bahais\ldots{} e assim continua
indefinidamente. Os Cingaleses e os Tâmeis\ldots{} Há sempre este confronto
entre grupos\ldots{} um tentando conquistar a terra ou o poder de outro.

Isso tem acontecido sabe-se lá desde quando. Houve sempre alguém a tentar
exterminar outra pessoa desde que Caim assassinou Abel -- e isso já foi há muito
tempo! Mas cada vez que ouvimos falar destas atrocidades, dizemos ``Que
terrível\ldots{} não devia acontecer\ldots{}''

Ouvimos falar de empresas farmacêuticas americanas que vendem drogas venenosas e
horríveis aos países do Terceiro Mundo. ``Isso não devia acontecer! É
terrível.'' A poluição do planeta, a espoliação do meio ambiente, a matança de
golfinhos e baleias\ldots{} quando é que isto acaba? Que podemos fazer por isto?
Parece ser um problema interminável de ignorância humana. Numa época em que as
pessoas deveriam saber melhor, elas estão a fazer as coisas mais horrendas umas
às outras. É uma época de previsões sombrias\ldots{} terremotos, erupções
vulcânicas e doenças\ldots{} não devia ser assim.

Agora dizer ``É assim que as coisas são'' não é uma aprovação, ou uma recusa
para se fazer qualquer coisa, mas é uma forma de nos estabelecermos no
conhecimento de que a Natureza é ``assim''. No reino animal a questão é muito
mais de sobrevivência do mais apto, uma lei natural de auto-selecção onde as
linhagens mais fracas são destruídas. Então, dessa forma, até a natureza é
bastante brutal, não é? Nós pensamos na Natureza como sendo tudo o que
``deveria'' ser\ldots{} doce, com flores e sol -- mas a Natureza também é muito
brutal.

Qual é a nossa posição na Natureza? Nós podemos viver no nível do reino animal,
com enfâse na sobrevivência do mais apto; o forte sobre o fraco e vivendo pelo
medo e poder. Podemos viver assim porque partilhamos essa mentalidade animal.
Temos um corpo animal que tem que sobreviver como qualquer outro corpo animal
neste planeta. Portanto, a Lei da selva é algo que os seres humanos podem
subscrever -- o que muitos deles fazem.

Mas este é apenas um nível inferior, certo? Se vivermos apenas nesse nível,
devemos esperar que o mundo seja como é -- num estado de medo e ansiedade. Mas,
como seres humanos, podemos ir além deste nível animal; podemos decidir nos
pautarmos por padrões morais, de modo a que não tenhamos que viver as nossas
vidas num estado de ansiedade.

Mas mais elevada ainda do que isso é a nossa capacidade de realizar a Verdade --
de contemplar a existência, de cultivar a mente reflexiva através da qual
podemos transcender a personalidade. No nível de comportamento moral, ainda
temos uma visão muito forte da personalidade. E ao longo da nossa civilização,
desenvolvemos um sentido de ``eu'' e ``meu'' a um nível absurdo. Tão forte é
esse sentido de ``eu'' e ``meu'' que parece dominar e contaminar tudo o que
fazemos e há sempre uma sensação de angústia e sofrimento ligada a isso.

Contemplemos apenas o seguinte: sempre que há uma sensação de ``eu'' e ``meu''
em qualquer coisa, isso parece gerar sempre descontentamento ou incerteza,
dúvida, culpa, medo ou ansiedade. Existe esta noção de ``eu'' como um ser
individual, de ``isto'' é meu, do que ``eu'' devia ou não devia, com base numa
crença de si mesmo como corpo ou condições mentais.

No entanto, essa visão é baseada em ilusão; vem do condicionamento, não de
compreensão interior. Portanto, desde que nos identifiquemos com as limitações
do corpo e da mente, é claro que vamos experimentar dúvida, desespero, angústia,
tristeza, mágoa e lamentação -- essas formas mentais de sofrimento. Como é que
poderia ser diferente? Nós não vamos certamente obter a iluminação com
mal-entendido distorcido e entendimento errado.

Temos esta oportunidade agora de estabelecer o Entendimento Correcto e a
Compreensão Correcta que nos liberta da ilusão da personalidade: a identificação
com o que é chamado os cinco aglomerados -- corpo, sensação, percepção,
formações mentais e consciência sensorial (cognição). Assim, contemplamos a
consciência sensorial através dos sentidos -- olhos, ouvidos, nariz, língua e
corpo. Podemos contemplar as formações mentais, as memórias dos passados da
nossa própria criação e os pensamentos e os pontos de vista que criamos. Podemos
vê-los como impermanentes.

Temos a capacidade de contemplar a natureza das coisas, de dizer: ``É assim que
as coisas são''. Podemos reparar `como as coisas são' sem adoptar um ponto de
vista de personalidade. Assim, com a respiração do corpo, o peso de sua postura,
estamos apenas a testemunhar e a observar como é \emph{agora}, neste momento. O
temperamento da mente, quer nos sintamos espertos ou entorpecidos, felizes ou
infelizes, é algo que podemos saber -- podemos testemunhar. E a mente vazia --
vazia das proliferações sobre si mesmo e os outros, é clareza. É inteligente e
compassiva. Quanto mais examinamos realmente os hábitos que desenvolvemos antes,
mais claras as coisas se tornam para nós. Portanto, devemos estar dispostos a
sofrer, ficar entediados, miseráveis e angustiados: é uma oportunidade de
suportar estes estados mentais desagradáveis, em vez de suprimi-los. Tendo
nascido, é o que é \emph{neste} momento, \emph{neste} lugar.
